\documentclass[10pt]{article}

% Specify the margins and text size.
\setlength{\textwidth}{6.4in}
\setlength{\textheight}{9.5in}
\setlength{\oddsidemargin}{0pt}
\setlength{\evensidemargin}{0pt}
\setlength{\topmargin}{0pt}
\setlength{\hoffset}{.05in}
\setlength{\voffset}{-1in}

\setlength{\parskip}{5pt}
\setlength{\parindent}{0pt}

% Load some fonts and symbol packages
\usepackage{latexsym}
\usepackage{pifont}       % contains 'star' symbol for counterinsurgency handbook title
\usepackage{yfonts} 
\usepackage{amsmath}
\usepackage{amsfonts}

\usepackage{graphicx}     % actually, this is loaded by pstricks
\usepackage[T1]{fontenc}
\usepackage{ifthen}
\usepackage{pstricks,pst-grad,pst-text,pst-node,multido,pst-plot,calc,pst-3dplot}
%\usepackage[all]{xy}
%\usepackage{animate}

% The hyperref package inserts links into the pdf.
\definecolor{MyLinkColor}{rgb}{.1,.2,1}
\definecolor{MyCiteColor}{rgb}{.1,1,.2}
\definecolor{MyURLColor}{rgb}{.4,.4,.4}
\usepackage[backref=true,pagebackref=false,hyperindex,colorlinks=true,
  linkcolor=MyLinkColor,urlcolor=MyURLColor]{hyperref}


% The tweaklist package is something I found on the web.  It provides a simple interface
% for making changes to spacing used in the itemize and enumerate environments.  Comment
% this out if you don't care to use tweaklist.
\usepackage{tweaklist}
\renewcommand{\itemhook}{\setlength{\parskip}{2pt}\setlength{\parsep}%
{1pt}\setlength{\topsep}{0pt}\setlength{\itemsep}{0pt}}

\newcommand{\U}{\underline{\hspace{5pt}}}

\usepackage{listings}
\newcounter{EX}\setcounter{EX}{1}
\newcommand{\EXERCISE}{\arabic{EX}.\stepcounter{EX} }

\begin{document}
\pagestyle{empty}
\lstset{language=R, showspaces=false, showstringspaces=false}

\href{http://www.su.edu}{\includegraphics[height=1.75cm]{sulogo.eps}}
\vspace{-1.69cm}

{\small{\ }\hfill
\begin{tabular}{cl}
& Math 207\\
& Introduction to Statistics\\
\end{tabular}
}
\setlength{\baselineskip}{1.05\baselineskip}
\bigskip

\begin{center}
\textbf{\large  The Binomial Formula}
\end{center}
\newcommand{\Z}{\hphantom{0}}

The chance that an event will occur exactly $k$ times out of $n$ is given by 
the binomial formula:
\[\frac{n!}{k!\,(n-k)!}\,p^k\,(1-p)^{n-k}\]
In this formula, $n$ is the number of trials, $k$ is the number of times
the event is to occur, and $p$ is the probability that the event will occur
on any particular trial.    The assumptions:
\begin{itemize}
\item The value of $n$ must be fixed in advance.
\item $p$ must be the same from trial to trial.
\item The trials must be independent.
\end{itemize}
\medskip

1. A coin is to be flipped 5 times. 
Use the binomial formula to calculate the following probabilities.

\hspace{20pt} a)   Exactly 3 heads
\vspace{.4in}

\hspace{20pt} b)   Exactly 4 heads
\vspace{.4in}

\hspace{20pt} c) Exactly 2 tails
\vspace{.4in}

\hspace{20pt} d) 3 or 4 tails
\vspace{.4in}

\hspace{20pt} e) At least 3 heads
\vspace{.4in}

2. Find the number of different ways of arranging one R and three Gs 
in a row.  Write out all the patterns.
\vspace{.4in}

3. Find the number of different ways of arranging two Rs and two Gs in a row.
Write out all the patterns.
\vspace{.4in}

4.  A die is rolled four times.  Find the chance that---

\hspace{20pt} a) an ace (one dot) never appears
\bigskip

\hspace{20pt} b) an ace appears exactly once
\bigskip

\hspace{20pt} c) an ace appears exactly twice.

\vfill
\eject
{\ }

5. A box contains one red marble and five green ones.  Four draws are made at 
random with replacement.  Find the chance that---

\hspace{20pt} a) a red marble is never drawn
\vspace{.4in}

\hspace{20pt} b) a red marble appears exactly once
\vspace{.4in}

\hspace{20pt} c) a red marble appears exactly twice
\vspace{.4in}

\hspace{20pt} d) a red marble appears exactly three times
\vspace{.4in}

\hspace{20pt} e) a red marble appears on all the draws
\vspace{.4in}

\hspace{20pt} f) a red marble appears at least twice
\vspace{.4in}
\vfill

5. (Review) An investigator wants to use a straight line to predict IQ from lead 
levels in the blood, for a representative group of children aged 5--9.  There is a weak
positive association in the data.  True or false, and explain---

\hspace{20pt} a) The investigator can use many different lines.
\bigskip
 
\hspace{20pt} b) The investigator has to use the regression line.
\bigskip

\hspace{20pt} c) Only the regression line has an r.m.s. error.
\bigskip

\hspace{20pt} d) Any line the investigator uses will have an r.m.s. error.
\bigskip

\hspace{20pt} e) Among all lines, the regression line has the smallest r.m.s. error.
\bigskip

\hspace{20pt} f) The r.m.s. error of his line will be $\mbox{SD}_y\,\sqrt{1-r^2}$.
\bigskip
%\vfill
\eject


\end{document}

